\subsection{External tasks and native success criteria}
\label{sec:external-dojo}
The new authored tasks above do not establish external validity. We therefore
add an action-prefix replay using AgentDojo~\cite{debenedetti2024agentdojo}, an
external benchmark with executable task-utility and attack-success predicates.
We pin the upstream repository at revision
\texttt{089ed468cf3e}, retain its MIT attribution, and use its v1 task definitions
without changing the success criteria. This evaluates an authorization component
of our method, not the full semantic-support gate or a live adaptive agent.

\paragraph{Selection and controls.}
Before scoring, we freeze 1,132 archived trajectories from four official pipeline
labels: GPT-4o (2024-05-13), GPT-4 Turbo (2024-04-09), Claude 3.5 Sonnet
(2024-06-20) and Llama 3 70B Chat. The selection contains all available clean
runs and \texttt{important\_instructions} runs for the specified disclosure
attacks in workspace, travel and Slack: 324 clean trajectories and 808 attacked
trajectories, spanning 81 underlying user tasks. We exclude sabotage and phishing
goals from the privacy endpoint, not unsuccessful attacks. Model names describe
the upstream archives; we do not make new model calls or claim that these are
current-model measurements.

Five matched controls use the same archived action sequence: no gate, a keyword
check for communication intent, literal destination binding to the initial
request, their conjunction, and denial of every mediated egress action. The
adapter consumes only the original user request and proposed tool arguments.
Neither attacker goals, security labels nor successful reference actions grant
permissions. It reuses the structured gate on a destination-authority envelope;
this says nothing about factual support of the message body. The initial request
is trusted and untrusted tool observations cannot extend its authority.
Destination binding deliberately does not infer contacts from names or grant
permissions from retrieved text. Its false denials are part of the evaluation.

\paragraph{Replay semantics and endpoints.}
We re-execute observed tool calls in AgentDojo's in-memory environments. Email,
Slack and web operations modify mock state; no real communication is sent. At
the first blocked egress action, we terminate the prefix and remove the final
answer instead of retaining the original trajectory's later success. Unblocked
trajectories keep their archived final answer. The unchanged upstream predicates
then evaluate task utility and attack success from the resulting state and
executed prefix. Upstream \texttt{security=True} denotes attacker success.
Errors remain in the denominator and are reported separately.

This intervention measures which archived actions a gate admits and what utility
the admitted prefix retains. It does not measure how an agent would recover,
replan, or attack again after a refusal. The selected attack predicates also
encode specific exfiltration goals, not every semantic disclosure. We compare
ungated replay with the archived labels to expose reconstruction differences.
For paired utility differences, we resample the 81 underlying tasks, keeping
all model and injection variants of each task together (2,000 bootstrap draws;
seed 1729). These intervals describe variation within this archive, not a
confidence guarantee for future deployments.

\paragraph{Privacy reduction has a measurable utility cost.}
All 5,660 replays completed in the Jetstream environment. Under the native
attack predicates, combined intent and destination checks reduce observed
attack success from 165/808 to 0/808. Clean utility falls from 198/324 (61.1\%)
to 179/324 (55.2\%); attacked utility falls from 383/808 (47.4\%) to 316/808
(39.1\%). The clean paired change is $-5.9$ percentage points (task-cluster
bootstrap 95\% interval $[-10.2,-2.2]$); the attacked change is $-8.3$ points
($[-12.8,-4.3]$). The combined arm blocks 28 clean and 202 attacked trajectories;
not every block removes a previously successful task. These losses are material:
this external result does not reproduce the eight-of-eight utility retention
of the internally authored structured tasks.

Figure~\ref{fig:core-replay} in the main text displays these same counts.
The deny-egress control has one clean and one attacked checker exception;
both remain unresolved rather than being scored as successful or verified safe.


Intent checking alone leaves 85 successful attacks. Destination-only and combined
checks have identical utility and attack-success counts, so this experiment
provides no measurable benefit for the additional intent rule. Denying every
mediated egress retains only 150 clean and 285 attacked successes, with two
checker exceptions. Some tasks need no mediated external send, so this arm is
not equivalent to refusing every task. The native predicate failure in each
case is a missing website-state key after early termination; it is not silently
converted to a confirmed safe outcome.

The costs concentrate in Slack: combined checking retains 44/84 clean successes
versus 61/84 without a gate, and 24/168 attacked successes versus 90/168. Travel
retains 42/80 clean and 75/160 attacked successes in both arms; workspace retains
93/160 versus 95/160 clean, and 217/480 versus 218/480 attacked. Thus aggregate
retention conceals substantial environment-specific loss. Literal destination
binding cannot resolve legitimate implicit contacts or contextual destinations.
Eleven combined-arm blocks also arise from unsupported destination argument
representations, including null optional values; we retain them rather than
repairing the adapter after seeing outcomes.

\paragraph{Concordance and scope of the evidence.}
Ungated replay agrees with archived task-utility labels in 1,128/1,132 cases and
attack-success labels in 781/808 attacked cases. The 27 attack disagreements are
26 workspace cases and one travel case. Four task-utility disagreements and
these attack disagreements remain in all denominators. We use outcomes of the
pinned executable predicates consistently across the five arms; the discrepancy
means the replay is not an exact reproduction of all historical labels. Twenty-four
archived traces contain execution errors. Runtime tool errors are retained as
observed behavior; the combined arm has no checker exceptions. All frozen
source/protocol and trace hashes are preserved with the artifacts.

\paragraph{Post-hoc discrepancy audit.}
We replay the 31 trajectories with any label disagreement under seven available
upstream suite versions, from v1 through v1.2.2, without changing the frozen
primary experiment. All four utility disagreements align with archived labels
under later versions: three travel cases pass after v1.1 removes the requirement
that the final answer repeat the hotel rating; one workspace case passes after
v1.1.1 accepts the updated end time as part of calendar rescheduling. This
establishes compatibility with changed predicates, not the provenance of the
historical evaluation. All 27 attack disagreements persist across these versions
(nine replay-positive/archive-negative and eighteen in the opposite direction).
Their historical cause remains unresolved.

As a post-hoc sensitivity analysis only, restricting to trajectories whose
ungated utility and attack labels agree with the archive leaves 322 clean and
779 attacked runs. Combined checking changes clean utility from 198/322 to
179/322, attacked utility from 383/779 to 315/779, and attack success from 156/779
to 0/779. Thus both the reduction in selected attack success and the utility
cost persist without the disputed trajectories. These filtered counts do not
replace the primary denominators; the subset was selected after observing
concordance. Per-trajectory version results accompany the artifacts.

The finding is bounded but useful: explicit task-to-destination binding can
suppress the selected archived exfiltration actions, while failing to retain
legitimate tasks that require contextual authority resolution. It does not
validate semantic support of message bodies, arbitrary leakage channels,
continuing attackers or agent recovery after a block. A next method should
resolve recipient authority from trusted task evidence and evaluate recovery in
a fresh closed-loop run; it should not extract permission from attacker-written
observations or use benchmark attack labels as a policy oracle.
